\section{The determinant identity for the terminal differential operators}
\label{rq-section}

The terminal matrix is not a matrix of commuting operators.
Its determinant identity has a spectral shift, and its
coefficients are functions which the derivatives can differentiate.
We prove the identity with both features retained.
The proof uses quantum minors, elementary exterior algebra,
and a polynomial cancellation argument.  It applies directly
to smooth functions and does not require an approximation by
regular functions.  The resulting identity is the source's
terminal determinant formula \cite{alpo-uniform-source}.

\subsection{Signed complementary minors}

Let $n\geq2$, let $\mathcal A$ be a unital complex algebra, and let
$E_{ab}\in\mathcal A$, $1\leq a,b\leq n$, satisfy
\[
 [E_{ab},E_{cd}]=\delta_{bc}E_{ad}-\delta_{da}E_{cb}.
\]
Use the ordered minors $Q_{I,J}(u)$ from
\eqref{rt-minor}.  Put
\[
 T_{ab}(u)=u\delta_{ab}+E_{ab},\qquad
 D(u)=Q_{\{1,\ldots,n\},\{1,\ldots,n\}}(u),
 \qquad C_{ij}(u)=(-1)^{i+j}Q_{\widehat i,\widehat j}(u).
\]
For distinct pairs define
\[
 N_{iqjp}(u)=
 \epsilon(i,q)\epsilon(j,p)(-1)^{i+q+j+p}
 Q_{\widehat{i,q},\widehat{j,p}}(u),
 \qquad
 \epsilon(i,q)=
 \begin{cases}1&i<q,\\-1&i>q.\end{cases}
\]
Set $N_{iqjp}=0$ when $i=q$ or $j=p$.
All omission lists are increasing.  Empty minors are one.
Here $C$ is indexed by the omitted row and omitted column;
it has not been transposed into an adjugate matrix.

We use the alternating and change-of-basis identities
for quantum minors proved in Theorem~\ref{rp-capelli}.
In particular $D(u)$ commutes with every $E_{ab}$:
the full determinant is unchanged by conjugation of the
generators, and differentiating this identity at the
identity matrix gives the asserted commutators.
It therefore commutes with every quantum minor.

\begin{lemma}\label{rq-cofactor-expansions}
For all indices and all $u\in\mathbb C$,
\begin{align}
 \sum_p C_{qp}(u)T_{cp}(u+1)
       &=\delta_{cq}D(u+1),                 \label{rq-cofactor-one}\\
 \sum_p N_{iqjp}(u)T_{cp}(u+1)
       &=\delta_{cq}C_{ij}(u+1)
                   -\delta_{ci}C_{qj}(u+1). \label{rq-cofactor-two}
\end{align}
\end{lemma}

\begin{proof}
The column expression for $Q$ and its reverse row
expression \eqref{rc-row-form} give the usual four
end-row and end-column expansions, with their displayed
shifts.  For example, expansion in the last row, written
with the last row factor on the right, gives
\[
 Q_{I,J}(v)=
 \sum_{\beta=1}^{s}(-1)^{s+\beta}
 Q_{I\setminus i_s,J\setminus j_\beta}(v-1)
 T_{i_sj_\beta}(v).
\]
An arbitrary row can be moved to the last position;
alternation contributes its permutation sign.
Equivalently, one may put an arbitrary row first in the
increasing-product row expression
\[
 Q_{I,J}(v)=
 \sum_{\tau\in S_s}\operatorname{sgn}(\tau)
 \prod_{\alpha=1}^{s}
 T_{i_\alpha,j_{\tau(\alpha)}}(v-s+\alpha).
\]
This expression follows by reversing the row list in
\eqref{rc-row-form} and using alternation.
These expansions are identities in $\mathcal A$;
they do not commute adjacent factors.

For \eqref{rq-cofactor-one}, append row $c$ to the list
with row $q$ removed, and expand along the appended row.
If $c\ne q$, a row is repeated and the minor is zero.
If $c=q$, moving that row back to its original position
gives $D(u+1)$ with the sign already included in
$C_{qp}$.

For \eqref{rq-cofactor-two}, start instead with the row
list omitting $i,q$ and the column list omitting $j$,
append row $c$, and expand in that row.
The result is zero unless $c=i$ or $c=q$.
In those two cases moving the appended row into its
increasing position gives respectively
$-C_{qj}(u+1)$ and $C_{ij}(u+1)$.
For precision, removing $q$ changes the position of
an index $i$ to $i-\mathbf1_{i>q}$, and removing $j$
changes the position of $p$ to $p-\mathbf1_{p>j}$.
The two resulting signs are exactly
$\epsilon(i,q)\epsilon(j,p)(-1)^{i+q+j+p}$.
This verifies every sign in
\eqref{rq-cofactor-two}, including the zero terms.
\end{proof}

\begin{theorem}\label{rq-jacobi}
The signed quantum minors satisfy
\begin{equation}\label{rq-jacobi-identity}
 C_{ij}(u+1)C_{qp}(u)-C_{qj}(u+1)C_{ip}(u)
                  =D(u+1)N_{iqjp}(u).
\end{equation}
No assumption that $\mathcal A$ is a domain is needed.
\end{theorem}

\begin{proof}
Fix $i,q,j$ and let $Z_p$ be the difference of the
two sides of \eqref{rq-jacobi-identity}.  Multiply
$Z_p$ on the right by $T_{cp}(u+1)$ and sum over $p$.
Lemma~\ref{rq-cofactor-expansions} gives
\begin{align*}
 \sum_p Z_pT_{cp}(u+1)
 &=\delta_{cq}C_{ij}(u+1)D(u+1)
       -\delta_{ci}C_{qj}(u+1)D(u+1)\\
 &\quad-D(u+1)
       \bigl(\delta_{cq}C_{ij}(u+1)
                       -\delta_{ci}C_{qj}(u+1)\bigr)=0.
\end{align*}
The last equality uses the centrality of $D$.
It remains to justify cancellation without assuming
invertibility of the matrix $T$.

Adjoin a central indeterminate $t$ and replace
$E_{ab}$ by $E_{ab}+t\delta_{ab}$.
These new generators satisfy the same Lie relations.
The preceding calculation gives polynomials
$Z_p(t)\in\mathcal A[t]$ satisfying
\[
 \sum_p Z_p(t)\bigl(E_{cp}+(u+1+t)\delta_{cp}\bigr)=0
                       \qquad\text{for every }c.
\]
If some $Z_p(t)$ is nonzero, let $d$ be the largest
degree of all of them.  The coefficient of $t^{d+1}$
in the equation with index $c$ is exactly the
degree-$d$ coefficient of $Z_c(t)$.  It vanishes
for every $c$, a contradiction.  Thus all $Z_p(t)$
are zero.  Substituting $t=0$ proves the theorem.
\end{proof}

\subsection{Moving the variable coefficients through a minor}

Suppose $\mathcal A$ also contains elements
$y_1,\ldots,y_n,\eta_1,\ldots,\eta_n$ such that
\begin{equation}\label{rq-vector-laws}
 [E_{ab},y_c]=-\delta_{ac}y_b,\qquad
 [E_{ab},\eta_c]=\delta_{bc}\eta_a.
\end{equation}
Assume that all the $y$'s and $\eta$'s commute with
one another.  They need not commute with the $E$'s.

\begin{lemma}\label{rq-moving}
Under these assumptions,
\begin{align}
 D(w)y_j
       &=y_jD(w)-\sum_l y_lC_{jl}(w),          \label{rq-det-vector}\\
 [C_{qp}(u),\eta_i]
       &=\sum_a N_{aqip}(u)\eta_a,             \label{rq-minor-covector}\\
 \sum_i\eta_iC_{ij}(u-1)
       &=\sum_iC_{ij}(u)\eta_i,               \label{rq-shift-one}\\
 \sum_a\eta_aN_{aqcp}(u-1)
       &=\sum_aN_{aqcp}(u)\eta_a.             \label{rq-shift-two}
\end{align}
\end{lemma}

\begin{proof}
For \eqref{rq-det-vector}, use the row expansion of
$D(w)$ which puts row $j$ in the far-left factor:
\[
 D(w)=\sum_l(-1)^{j+l}
                 T_{jl}(w-n+1)Q_{\widehat j,\widehat l}(w).
\]
Every generator in the smaller minor commutes with
$y_j$, because its row is different from $j$.
Commuting $y_j$ through the remaining factor gives
$[T_{jl}(w-n+1),y_j]=-y_l$.
This proves \eqref{rq-det-vector}.

For \eqref{rq-minor-covector}, the assertion is zero
if $i=p$, since column $p$ is absent.
Otherwise expand the minor defining $C_{qp}(u)$
in column $i$, putting that column factor on the right.
All factors in the smaller minors commute with
$\eta_i$ because their columns differ from $i$.
The remaining commutator is
$[T_{ai}(u-n+2),\eta_i]=\eta_a$.
The position of $a$ after removing $q$ is
$a-\mathbf1_{a>q}$ and that of $i$ after removing
$p$ is $i-\mathbf1_{i>p}$.
The signs from these positions, together with
$(-1)^{q+p}$, are precisely those in
$N_{aqip}(u)$.  This proves the second identity.

Here is an exterior-algebra proof of both shift
identities, including their parameter change.
Adjoin symbols $\theta_1,\ldots,\theta_n$ commuting
with $\mathcal A$ and satisfying
$\theta_a\theta_b=-\theta_b\theta_a$ and $\theta_a^2=0$.
Set
\[
 \beta_j(v)=\sum_a\theta_aT_{aj}(v),\qquad
                    \eta=\sum_a\theta_a\eta_a.
\]
Relation \eqref{rq-vector-laws} gives, term by term,
\[
                   \beta_j(v)\eta=-\eta\beta_j(v-1).
\]
Indeed the commutator contribution is
$\sum_a\theta_a\theta_j\eta_a=-\theta_j\eta$.
Repeatedly using this equality yields
\[
 \beta_{j_1}(u)\beta_{j_2}(u-1)\cdots
       \beta_{j_s}(u-s+1)\eta
   =(-1)^s\eta\,\beta_{j_1}(u-1)\cdots\beta_{j_s}(u-s).
\]
The coefficient of
$\theta_{i_1}\cdots\theta_{i_{s+1}}$ is
\begin{equation}\label{rq-general-shift}
 \sum_{\alpha=1}^{s+1}(-1)^{\alpha-1}
       \eta_{i_\alpha}Q_{I\setminus i_\alpha,J}(u-1)
  =
 \sum_{\alpha=1}^{s+1}(-1)^{\alpha-1}
       Q_{I\setminus i_\alpha,J}(u)\eta_{i_\alpha}.
\end{equation}
This follows directly from the column definition
of the minors.  Taking $I$ to be the full index
list and $J$ to omit $j$ proves
\eqref{rq-shift-one}.  Taking $I$ to omit $q$
and $J$ to omit $c,p$, then multiplying by the
fixed omission signs, proves \eqref{rq-shift-two}.
\end{proof}

\subsection{The contracted determinant identity}

For fixed complex vectors $x=(x_j)$ and
$\xi=(\xi_i)$ define
\begin{align*}
 S(u;x,\xi)&=\sum_{i,j}\xi_i x_j C_{ij}(u),\\
 H_y(u;\xi)&=\sum_{i,j}\xi_i y_j C_{ij}(u),\\
 H_\eta(u;x)&=\sum_{i,j}x_j\eta_i C_{ij}(u),\\
 H_{y\eta}(u)&=\sum_{i,j}\eta_i y_j C_{ij}(u),\\
 H_2(u;x,\xi)&=
 \sum_{\substack{i<i'\\j<j'}}(-1)^{i+i'+j+j'}
  (\eta_i\xi_{i'}-\eta_{i'}\xi_i)
  (y_jx_{j'}-y_{j'}x_j)
            Q_{\widehat{i,i'},\widehat{j,j'}}(u).
\end{align*}
In every expression the coefficient precedes the minor.

\begin{theorem}\label{rq-contracted}
One has the identity in $\mathcal A$
\begin{equation}\label{rq-contracted-identity}
 D(u+1)H_2(u-1;x,\xi)
     +H_y(u+1;\xi)H_\eta(u-1;x)
                    =H_{y\eta}(u)S(u;x,\xi).
\end{equation}
\end{theorem}

\begin{proof}
Both sides are bilinear in $x,\xi$, so fix $p,q$
and take $x_j=\delta_{jp}$, $\xi_i=\delta_{iq}$.
For brevity write
\[
 C^+_{ij}=C_{ij}(u+1),\quad
 C^0_{ij}=C_{ij}(u),\quad C^-_{ij}=C_{ij}(u-1),
 \quad N^0_{aqcp}=N_{aqcp}(u).
\]
Expanding the two alternating factors in $H_2$
and grouping their four terms gives
\[
 H_2(u-1;x,\xi)=\sum_{a,c}\eta_a y_cN_{aqcp}(u-1).
\]
By \eqref{rq-shift-two} and the commutativity of
$\eta_a$ with $y_c$, this is
$\sum_c y_c\sum_aN^0_{aqcp}\eta_a$.
Likewise \eqref{rq-shift-one} gives
\[
 H_\eta(u-1;x)=\sum_iC^0_{ip}\eta_i,\qquad
 H_{y\eta}(u)=\sum_l y_l\sum_i C^+_{il}\eta_i.
\]
Let
\[
                  \mathcal E=
       \sum_{l,c,a}y_l C^+_{cl}N^0_{aqcp}\eta_a .
\]
Using \eqref{rq-det-vector} in the first summand
of the proposed identity gives
\[
 D(u+1)H_2(u-1;x,\xi)
       =\sum_{l,a}y_lD(u+1)N^0_{aqlp}\eta_a-\mathcal E.
\]
The second summand is
$\sum_{l,a}y_l C^+_{ql}C^0_{ap}\eta_a$.
Their sum is therefore
\[
 \sum_{l,a}y_l
   \bigl(D(u+1)N^0_{aqlp}+C^+_{ql}C^0_{ap}\bigr)\eta_a
                       -\mathcal E
   =\sum_{l,a}y_l C^+_{al}C^0_{qp}\eta_a-\mathcal E
\]
by Theorem~\ref{rq-jacobi}.
On the other hand,
\eqref{rq-minor-covector} gives
\[
 H_{y\eta}(u)C^0_{qp}
       =\sum_{l,a}y_l C^+_{al}C^0_{qp}\eta_a-\mathcal E.
\]
This is the same expression.  Bilinearity now proves
\eqref{rq-contracted-identity}.
\end{proof}

\subsection{Application to the actual terminal group}

Return to $W$ of signature $(k,1)$ and its group
$U(W,h)$, with $h$ linear in the first argument.
For a smooth function on this group let $E_{ab}$
be its complexified right derivatives.
Multiplication by
\[
 y_j(m)=(m^{-1}z)_j,\qquad
       \eta_i(m)=h(mw_i,z)=h(w_i,m^{-1}z)
\]
defines commuting operators.
Differentiation of $m^{-1}z$ and of the row
$z^*Hm$ gives exactly
\eqref{rq-vector-laws}.
Take $x_j=z'_j$, $\xi_i=h(w_i,z')$ and
$u=Y+k/2$.  Equations
\eqref{rc-t21}--\eqref{rc-t22} become
\[
 t_{11}(Y)=D(u),\quad
 t_{21}(Y)=-H_y(u;\xi),\quad
 t_{12}(Y)=-H_\eta(u-1;x),\quad
 t_{22}(Y)=-H_2(u-1;x,\xi).
\]
Thus Theorem~\ref{rq-contracted} proves
\begin{equation}\label{rq-terminal-factorization}
 t_{11}(Y+1)t_{22}(Y)-t_{21}(Y+1)t_{12}(Y)
                         =-H_{y\eta}(u)S(u;x,\xi).
\end{equation}
All compositions here are compositions of genuine
differential operators with variable coefficients.

Choose orthonormal bases $B_1,\ldots,B_k$ of $z^\perp$
and $B'_1,\ldots,B'_k$ of $(z')^\perp$ in $W$.
Write $C^R_{B'}(u)$ for the degree-$k$ Capelli
determinant of the right derivatives along
$B'_a h(\,\cdot\,,B'_b)$.
Write $C^L_B(u)$ for the same polynomial in left
derivatives, with every product of derivatives reversed.
This reversal is part of the definition.

\begin{theorem}\label{rq-terminal-determinant}
For every smooth function $F$ on the terminal group,
\begin{equation}\label{rq-exact-terminal}
 \bigl(t_{11}(Y+1)t_{22}(Y)-t_{21}(Y+1)t_{12}(Y)\bigr)F
          =-TT'\,C^L_B(Y+k/2)C^R_{B'}(Y+k/2)F.
\end{equation}
The formula holds for every complex $Y$ and needs
neither nonvanishing intermediate minors nor a
positivity restriction on the Gram polynomial.
\end{theorem}

\begin{proof}
In the ordered basis
$B'_1,\ldots,B'_k,z'/\sqrt{T'}$,
the vector $x$ has only its last entry nonzero,
equal to $\sqrt{T'}$, and the covector $\xi$
has last entry $-\sqrt{T'}$.
The change-of-basis rule for minors therefore gives
\[
                       S(u;x,\xi)=-T'C^R_{B'}(u).
\]
At a fixed $m$, apply the same calculation to
$m^{-1}z$, whose Hermitian norm is $-T$.
An orthonormal frame of its perpendicular space is
$m^{-1}B_1,\ldots,m^{-1}B_k$.  Since the
coefficients in $H_{y\eta}$ precede the derivatives,
they are frozen at this evaluation point.  Consequently
\[
 (H_{y\eta}(u)G)(m)
                    =-T(C^R_{m^{-1}B}(u)G)(m).
\]
There is no derivative of the frozen frame in this
identity.

For completeness, if $X_a=\operatorname{Ad}(m^{-1})Y_a$
are fixed tangent matrices at this point, then
\[
 (R_{X_1}\cdots R_{X_s}G)(m)
                       =(L_{Y_s}\cdots L_{Y_1}G)(m).
\]
Indeed both sides are the mixed derivative at zero
of
$G(\exp(t_1Y_1)\cdots\exp(t_sY_s)m)$.
Complexification gives the same identity for the
matrix units defining the Capelli determinant.
This proves
$(C^R_{m^{-1}B}(u)G)(m)=(C^L_B(u)G)(m)$,
including the required reversal of each product.
Substitution into
\eqref{rq-terminal-factorization} proves
\eqref{rq-exact-terminal}.
\end{proof}

\begin{corollary}\label{rq-type-scalar}
Suppose
$F(hmh')=\pi(h)F(m)\pi'(h')$ for
$h\in U(z^\perp)$ and $h'\in U((z')^\perp)$,
where $\pi,\pi'$ are irreducible finite-dimensional
unitary representations with highest weights
$\lambda=(\lambda_1\geq\cdots\geq\lambda_k)$ and
$\lambda'=(\lambda'_1\geq\cdots\geq\lambda'_k)$.
Then the left side of \eqref{rq-exact-terminal} equals
\[
 -TT'\prod_{j=1}^k
   (Y+k/2-j+1+\lambda_j)
   (Y+k/2-j+1+\lambda'_j)\,F.
\]
In the source's parameters
$u_j=\lambda_j+(k+1)/2-j$ and
$u'_j=\lambda'_j+(k+1)/2-j$, this scalar is
\[
                 -TT'\prod_{j=1}^k
                    (Y+\tfrac12+u_j)(Y+\tfrac12+u'_j).
\]
\end{corollary}

\begin{proof}
The right derivatives act through the derived
representation $\pi'$ on the right of $F$.
Left derivatives act in the reverse order through
$\pi$ on its left; the reversal in $C^L_B$
therefore restores the ordinary Capelli polynomial.
Its value in highest weight $\lambda$ is
$\prod_j(u-j+1+\lambda_j)$ by
Theorem~\ref{rp-capelli}.
Apply this on both sides of
\eqref{rq-exact-terminal} and substitute $u=Y+k/2$.
\end{proof}

\paragraph{Scope.}
This proves the terminal determinant formula and its
actual finite-type scalar, including all variable
coefficient corrections.  It does not assert the
remaining commutation relations of the complete
terminal matrix, the recursive frame identity, or
integration by parts for the Jacquet pairing.
Those assertions have their own hypotheses and are
not consequences of this determinant identity alone.

The spectral shift cannot be suppressed.  In dimension
two, take the trivial representation $E_{ab}=0$ and
$u=1$.  Then $C_{11}(1)C_{22}(1)=1$, while
$D(1)N_{1212}(1)=0$.  The unshifted version of
\eqref{rq-jacobi-identity} is therefore false even in
this elementary representation.  The submitted
identity has the shift proved above, so this is a
counterexample only to that stronger, altered statement.


